\section{Semiample extensions of the torus and symplectic factors}
\label{periods:section}

We now treat the nonprojective factors in Proposition~\ref{spread:product}.
Their first or second cohomology has a classical period domain.  Two points
require care: the available K\"ahler classes are real, and the compact auxiliary
base need not be algebraic.  We first address the polarization and extension
questions independently of the factor construction.

\subsection{Real polarizations and the integral local system}

An integral variation in this section consists of a local system
$\mathbb V_{\Z}$ of finite free abelian groups and a holomorphic Hodge filtration
on $E=\mathbb V_{\Z}\otimes_{\Z}\cO$, satisfying opposedness and Griffiths
transversality.  A real polarization is a flat nondegenerate real bilinear
form on $\mathbb V_{\R}$ satisfying the Hodge--Riemann relations.  We use the
sign convention in which, on type $(p,q)$ of weight $w$, the associated
positive Hermitian form is
\begin{equation}\label{periods:sign}
 (-1)^{w(w-1)/2}(\sqrt{-1})^{p-q}B(v,\overline v).
\end{equation}
Changing this convention by the usual weight-dependent sign changes none of
the arguments.

\begin{lemma}[Changing the polarization]\label{periods:rationalization}
Let $(\mathbb V_{\Z},F^\bullet)$ be an integral pure variation with a flat real
polarization $B$.  There is a flat real automorphism $C$ of $\mathbb V_{\R}$
such that $(\mathbb V_{\Z},C^{-1}F^\bullet)$ is polarized by a rational form
$B'$, and
\[
 B'(x,y)=B(Cx,Cy).
\]
Consequently the new and original variations are isomorphic as real
polarized variations, and their holomorphic Hodge bundles are isomorphic.
Their canonical Hodge extensions are also identified wherever local
monodromy is unipotent.  After multiplying $B'$ by a positive integer, its
values on $\mathbb V_{\Z}$ are integral.
\end{lemma}

\begin{proof}
Work on a connected base and fix a lattice fiber $\Lambda$ and its monodromy
representation.  In a basis of $\Lambda$, the invariant bilinear forms of the
required symmetry satisfy rational linear equations
\[
 g^{\mathsf t}Bg=B\quad(g\text{ in the monodromy group}),
 \qquad B^{\mathsf t}=(-1)^wB.
\]
Although the first collection may be infinite, the rows of these equations
span a finite-dimensional rational vector space.  Finitely many suffice.
Thus rational invariant forms are dense in the space of real invariant
forms.  Choose such a form $B'$ sufficiently close to $B$.

Put $J=B^{-1}B'$.  Both invariance identities imply that $J$ commutes with
monodromy, and the symmetry identities give
$J^{\mathsf t}B=BJ$.  If $B'$ is sufficiently close, the power series for the
square root at the identity defines $C=J^{1/2}$.  This is a real invertible
matrix, commutes with monodromy, and satisfies $C^{\mathsf t}B=BC$.
Therefore
\[
 C^{\mathsf t}BC=BC^2=BJ=B'.
\]
In particular $C$ defines a global flat real automorphism.  The filtration
$F'^p=C^{-1}F^p$ is holomorphic and transverse because $C$ is flat; opposedness
holds because $C$ is real.  The displayed identity transfers the
Hodge--Riemann relations exactly, at every base point.  Hence $B'$ polarizes
$F'^\bullet$.

The lattice and its monodromy have not changed.  We do not require $C$ to
preserve the lattice: it supplies an isomorphism of real variations and of
their holomorphic filtered bundles.  On a unipotent normal-crossing chart,
$C$ commutes with every monodromy logarithm.  In the canonical logarithmic
frames it is therefore the same invertible constant matrix.  It identifies
the flat extensions and the extended Hodge filtrations.  Finally, clear the
denominators of $B'$ on a lattice fiber; monodromy invariance clears them
everywhere.  If this multiplies $B'$ by $a>0$, replace $C$ by $\sqrt a\,C$;
the transported filtration is unchanged and the polarization identity
remains exact.
\end{proof}

Thus Lemma~\ref{periods:rationalization} does not assert that a nonprojective
fiber has a rational polarization of its original Hodge structure.  It
constructs a different Hodge filtration on the same integral local system,
with isomorphic holomorphic Hodge bundles.  This is precisely what is needed
to use arithmetic period spaces for these bundles.

\begin{lemma}[Polarizations from a fixed K\"ahler class]
\label{periods:lefschetz}
Let $t:\mathcal T\to B^\circ$ be a smooth proper family of connected compact
complex manifolds of dimension $n$ over a complex manifold $B^\circ$,
and suppose a closed real $(1,1)$-form
on $\mathcal T$ restricts to a K\"ahler form on every fiber.  Its restriction
class $\kappa$ is flat.  The cohomological variations, with their integral
lattices modulo torsion, are real-polarizable.  In particular:
\begin{enumerate}[label=\textup{(\roman*)}]
\item on $H^1$, a polarization is
\[
 B_1(\alpha,\beta)=\int_{\mathcal T_b}
                  \alpha\wedge\beta\wedge\kappa^{n-1};
\]
\item if $n\geq2$, write
$H^2_{\R}=P^2_{\R}\oplus\R\kappa$, where
$P^2_{\R}=\Ker(\kappa^{n-1}\cup-)$.
For $\alpha=\alpha_0+a\kappa$ and $\beta=\beta_0+b\kappa$, a full
weight-two polarization is
\begin{equation}\label{periods:full-h2}
 B_2(\alpha,\beta)=
 \int_{\mathcal T_b}\alpha_0\wedge\beta_0\wedge\kappa^{n-2}
 -ab\int_{\mathcal T_b}\kappa^n.
\end{equation}
If $h^{2,0}=1$, this real form has signature $(2,h^{1,1})$.
\end{enumerate}
\end{lemma}

\begin{proof}
Restriction of a closed total-space class is invariant under parallel
transport in a smooth proper family.  For example, a differentiable
trivialization over a path and Stokes' theorem show that its integrals on
transported cycles are constant.  Thus $\kappa$, its cup powers, and the
fiber integrals in the statement are flat.  Locally on the base, properness
allows us to add a sufficiently large K\"ahler form pulled back from a
coordinate ball and obtain a K\"ahler form on the total space.  The smooth
proper K\"ahler Hodge
theorem gives the Hodge filtrations and transversality; one may use the
smooth, boundary-free case of \cite[Theorem~1.1]{FF25}.  The K\"ahler
Lefschetz decomposition and Hodge--Riemann relations then polarize its
primitive summands.  Transporting their forms by Lefschetz and inserting
the signs for the resulting Tate twists polarizes full cohomology.

For clarity, $H^1$ is primitive and its Hodge--Riemann form is $B_1$.
In degree two the summands $P^2_{\R}$ and $\R\kappa$ are orthogonal for the
intersection form with $\kappa^{n-2}$.  This form polarizes the primitive
summand with convention~\eqref{periods:sign}.  The positive K\"ahler line
must have its sign reversed, exactly as in~\eqref{periods:full-h2}.
The resulting form is positive on the real plane underlying $H^{2,0}$
and negative on real $H^{1,1}$ when $h^{2,0}=1$.  This proves the asserted
signature.  The primitive projectors and these forms are flat, so all
constructions apply to variations.
\end{proof}

For the families in Proposition~\ref{spread:product}, the form in this
Lemma is obtained by pulling back a fixed K\"ahler form from the compact
ambient manifold $V$.  Locally over a small polydisk in the auxiliary base,
the family embeds in the product of that polydisk with $V$, so its total
space is K\"ahler there.  No K\"ahler metric on the entire auxiliary base
is required.

\subsection{Canonical extensions and volume norms}

For unipotent monodromy along a simple normal crossing divisor, we write
$\overline E$ for the Deligne extension with nilpotent residues, and
$\overline F^p$ for its Schmid Hodge subbundles.  Their existence in the
analytic setting follows from \cite[Theorem~4.12]{Schmid73}; this theorem
concerns a product of punctured disks and disks.  By
Lemma~\ref{periods:rationalization}, its usual polarized integral
formulation applies equally to the real-polarized variations used here.

\begin{lemma}[Functoriality and logarithmic bounds]
\label{periods:extensions}
Let a pure real-polarizable integral variation be defined on the complement
of an SNC divisor in a complex manifold, with unipotent local monodromy.
Canonical Hodge extensions commute with tensor products, exterior powers,
duals, and holomorphic pullbacks from smooth manifolds with SNC boundary,
whose boundary complements map into the original complement.  This
includes disks whose punctured disks map into that complement.
A morphism of pure real variations induces an isomorphism
between its coimage and image on these extensions.  In particular, an
isomorphism of highest Hodge lines induced by such a morphism remains an
isomorphism after extension.

If $\mathcal H$ is a highest Hodge line, $\gamma:\Delta\to B$ is a disk
with $\gamma(\Delta^*)$ in the open set, and $e$ is a nowhere-zero local
section of $\gamma^*\overline{\mathcal H}$, there are constants $c>0$ and
$a\geq0$ such that, for $0<|z|$ sufficiently small,
\begin{equation}\label{periods:logbounds}
 c^{-1}(1+|\log|z||)^{-a}
 \ \leq\ \|e(z)\|\ \leq\
 c(1+|\log|z||)^a.
\end{equation}
The same assertion holds for its powers and their duals.
\end{lemma}

\begin{proof}
On a normal-crossing chart let $N_j$ be the commuting logarithms of the
local monodromies.  Applying the exponential of
$-\sum_j(\log z_j)N_j/(2\pi\sqrt{-1})$ to multivalued flat frames gives
canonical logarithmic frames, with a compatible choice of monodromy
convention.  Tensor and dual constructions have the corresponding sums
and duals of the nilpotent residues, so their flat extensions agree.
Their extended Hodge filtrations agree as well: the induced filtrations
are holomorphic subbundles, and coincide on the dense open with the given
ones.

After a holomorphic pullback the residue along a prime divisor is a sum
$\sum_j a_jN_j$, where the $a_j$ are the orders of the pulled-back boundary
coordinates.  It is nilpotent because the $N_j$ commute.  Thus the pulled-back
flat extension is again canonical; pulling back the Hodge subbundles gives
their canonical extension by the same uniqueness.  Locally on a disk,
write the coordinates as $z^{a_j}u_j(z)$ with nonvanishing holomorphic
units $u_j$; choosing logarithms of the units gives the explicit
identification.  This also explains the usual pullback compatibility
\cite[Lemma~5.6]{BFMT25} in the analytic situation needed here.

For a morphism of pure variations, Lemma~\ref{hodge:purerules} supplies
flat Hodge splittings of its kernel and image.  In logarithmic frames
these splitting projectors extend as the same constant matrices, so the
splittings persist in the extended filtrations.  This proves the
coimage-to-image assertion and its consequence for highest lines.

It remains to check the norm statement.  Pull back to the disk, using the
compatibility just proved.  Schmid's abstract norm estimate
\cite[Theorem~6.6$'$]{Schmid73}, applied to a flat basis and to the dual
variation, bounds both the Hodge metric and its inverse by powers of
$1+|\log|z||$ on sectors.  The matrices relating a flat basis to a
canonical logarithmic basis are polynomials in $\log z$; hence the same
kind of estimates holds in canonical frames.  Finitely many sectors
cover a punctured disk.  The coefficients of $e$ in a holomorphic frame
of $\overline E$ are bounded.  Since its highest line is a subbundle and
$e$ does not vanish, a local holomorphic section $\xi$ of
$\overline E^{\vee}$ can be chosen with $\xi(e)=1$.  Its coefficients are
also bounded.  The metric estimates and
$1\leq\|\xi\|\,\|e\|$ give both inequalities in
\eqref{periods:logbounds}.  Powers and duals follow immediately.
\end{proof}

In degree $n$ on an $n$-dimensional compact K\"ahler manifold, a
holomorphic $n$-form is primitive.  Its Hodge norm is, with a fixed
normalization,
\begin{equation}\label{periods:volume}
 \|\eta\|_{\mathrm{vol}}^2
   =(\sqrt{-1})^{n^2}\int_{\mathcal T_b}\eta\wedge\overline\eta.
\end{equation}
Thus Lemma~\ref{periods:extensions} gives bounds for the actual volume norms
appearing in the product formula, as well as for abstract Hodge norms.

\subsection{Arithmetic period maps on an analytic base}

For an effective pure variation of weight $w$, we use
the Griffiths line
\[
 \lambda=\bigotimes_{p=1}^{w}\det F^pE.
\]
Including the full step $F^0E=E$ only adds a full-local-system determinant;
we will account for this harmless finite-order factor explicitly.

We use the following precise consequence of
\cite[Theorems~5.2 and~5.5]{BFMT25}.  If $Z$ is an algebraic variety with
a quasifinite horizontal period map for a polarized integral variation,
then it has a projective period compactification $Z^{\mathrm{BB}}$.
For a sufficiently divisible positive integer $a$ an ample line bundle
$\mathcal A$ on this compactification restricts to $\lambda_Z^a$.
Every analytic map $(\Delta^*)^b\to Z^{\mathrm{an}}$ whose period map is
locally liftable extends to $\Delta^b\to Z^{\mathrm{BB},\mathrm{an}}$;
when the induced local monodromies are unipotent, the pullback of
$\mathcal A$ is canonically the Schmid extension of the corresponding
$a$-th Griffiths power.  The last assertion is part of Theorem~5.5, not an
algebraicity assumption on the source.

\begin{lemma}[The Griffiths line on a compact analytic base]
\label{periods:griffiths}
Let $B$ be a smooth compact complex manifold, let $D$ be an SNC divisor,
and let $\mathbb V$ be a real-polarizable integral variation on
$B^\circ=B\setminus D$.  Suppose its Hodge numbers are either
\[
  w=1,\quad h^{1,0}=h^{0,1}=g,
  \qquad\text{or}\qquad
  w=2,\quad h^{2,0}=h^{0,2}=1.
\]
After a finite cover of $B^\circ$, extended and resolved over $B$, the
canonical extension of its Griffiths line is semiample.  The resulting
base is smooth and compact and its boundary can be taken SNC.
\end{lemma}

\begin{proof}
Apply Lemma~\ref{periods:rationalization} and scale the rational
polarization to be integral.  In weight one its domain $\mathcal D$ is
Siegel space.  In weight two, use the sign convention giving signature
$(2,h^{1,1})$.  The highest filtration step is a line $\ell$ satisfying
$B'(\ell,\ell)=0$ and $B'(v,\overline v)>0$ for $v\ne0$ in $\ell$;
the middle step is $\ell^\perp$.  A connected component of this domain
is a type-IV domain.  A zero-dimensional domain causes no difficulty and
can be treated as a point.

Let $\Gamma$ be the integral isometry group preserving a component, after
passing to the corresponding finite cover if necessary.  It is an
arithmetic group and contains the monodromy.  Choose a finite-index neat
subgroup $\Gamma'$, as in the level construction of
\cite[\S5.2]{BFMT25}.  Here neat means that the multiplicative group
generated by the eigenvalues of each element has no nontrivial torsion.
The inverse image of $\Gamma'$ under monodromy gives a finite \emph{\'etale}
cover of $B^\circ$.  This cover has a finite normal extension over $B$
without any K\"ahler assumption on $B$.  Indeed, on an SNC chart
$\Delta^{*k}\times\Delta^{b-k}$ its permutation monodromy is finite.
A sufficiently divisible coordinate power map kills that monodromy.
The cover is then a quotient of finitely many copies of this power
cover, with its finite deck group also permuting the copies.  Extend
each copy to the full polydisk and take the same finite-group quotient.
This is a finite normal analytic extension.  Such extensions are unique,
being the normalization in the finite meromorphic algebra defined on the
open, so they glue.  The resulting normal space is compact because it
is finite over $B$.  Resolve it with SNC boundary, leaving the smooth
covering open unchanged.  This is the finite-extension construction also
used in the proof of \cite[Lemma~2.21]{Villadsen24}; the additional
K\"ahler conclusion of that Lemma is not needed here.

Every boundary monodromy on this resolution lies in $\Gamma'$.
By the monodromy theorem, in its analytic formulation
\cite[Lemma~4.5]{Schmid73}, all its eigenvalues are roots of unity.
Neatness therefore makes them all equal to one.  Hence all local
monodromies are unipotent, including those about exceptional boundary
divisors.  We keep the notation $B$ for the resulting base.

The quotient $Z=\Gamma'\backslash\mathcal D$ is a smooth algebraic
quasiprojective variety by Baily--Borel
\cite[\S10, especially Theorem~10.11]{BailyBorel66}.  The theorem applies
to the effective adjoint group of the classical domain, and the finite
central quotient does not change this conclusion.  We keep $\Gamma'$
inside the original integral isometry group: neatness kills the finite
kernel of its action and the finite point stabilizers.  Thus the local
system $(\mathcal D\times V_{\Z})/\Gamma'$, with $V_{\Z}$ the marked
lattice, and the tautological
filtration give an integral polarized variation on $Z$, retaining the
original representation.  In the zero-dimensional weight-two case,
the isometry group of the positive-definite integral lattice is finite,
so its neat subgroup is trivial and this variation is constant on a
point.  Its period map to this
arithmetic quotient is the identity, so it is quasifinite.  It is
horizontal: for weight one transversality is automatic, and in weight
two differentiation of $B'(v,v)=0$ gives $dv\in\ell^\perp$.
Thus $Z$ satisfies the Griffiths-line hypotheses of
\cite[Theorems~5.2 and~5.5]{BFMT25}.

The conjugated variation on $B^\circ$ is the pullback of this tautological
variation by its analytic period map $\phi$.  This assertion concerns
the local system as well as the filtration: the marked period map is
equivariant for its monodromy representation, and at neat level the
quotient has no stabilizers.  In particular $\phi$ is locally liftable.
The quoted analytic extension theorem extends it to
\[
 \overline\phi:B\longrightarrow Z^{\mathrm{BB},\mathrm{an}},
 \qquad
 \overline\phi^*\mathcal A\simeq\overline\lambda^{\,a}.
\]
To apply the theorem on a chart $\Delta^{*k}\times\Delta^{b-k}$,
restrict first to $(\Delta^*)^b$.  Its extension agrees with the original
map on $\Delta^{*k}\times\Delta^{b-k}$ by uniqueness.  The unused
coordinates have trivial monodromy, so the same statement holds for the
canonical bundle identification.  Both maps and identifications glue
uniquely on overlapping charts.

A sufficiently high power of $\mathcal A$ is generated by global
sections on the projective variety $Z^{\mathrm{BB}}$.  Pulling these
sections back gives global generation of a power of $\overline\lambda$
on the compact analytic base.  Finally, the real isomorphism in
Lemma~\ref{periods:rationalization} identifies this extension with the
one for the original variation.
\end{proof}

\subsection{From Griffiths determinants to the required highest lines}

We record the determinant calculation as an identity of holomorphic
bundles.  For any integral local system, $\det E$ has monodromy in
$\{1,-1\}$, since its monodromy matrices are invertible over $\Z$.
If the boundary monodromies are unipotent, its canonical extension has
trivial local monodromy and extends as a flat line across the boundary.
Its square is therefore holomorphically trivial on the compactification.
At neat level the determinant itself is trivial.  These statements concern
the full determinant, and do not assert finite monodromy of individual
Hodge pieces.

In weight one, the Griffiths line is $\det F^1E$ up to this full
determinant.  In weight two put $\ell=F^2E$.  The polarization identifies
$E/F^1E$ with $\ell^\vee$, giving
\begin{equation}\label{periods:determinants}
 \det F^1E\simeq\det E\otimes\ell,
 \qquad
 \lambda\simeq\det E\otimes\ell^{\otimes2}.
\end{equation}
The possible additional $F^0$ factor changes only the power of $\det E$.
These identities hold on canonical extensions.  Indeed, the flat
polarization extends as a nondegenerate bilinear form in logarithmic
frames, and $\overline F^1=(\overline F^2)^\perp$ by continuity of the
subbundles.  Taking determinants then gives the extended identities.
It follows that semiampleness of $\overline\lambda$ implies semiampleness
of $\overline\ell$: kill the determinant torsion and use a further even
power.  Thus the square in~\eqref{periods:determinants} is retained.

\begin{proposition}[Semiample period-factor extensions]
\label{periods:semiample}
Let $r:P\to S$ and the smooth families $t_i:T_i\to P^\circ$ be as in
Proposition~\ref{spread:product}, with $n_i=\dim T_{i,u}>0$, and put
\[
 \mathcal H_i=F^{n_i}\bigl(R^{n_i}(t_i)_*\C\otimes\cO_{P^\circ}\bigr).
\]
Thus $\mathcal H_i$ is the highest line denoted $\lambda_i$ in
Proposition~\ref{spread:product}; the symbol $\lambda$ above denotes
the Griffiths line.
After a proper generically finite surjection $\pi:P'\to P$, obtained by
extending a finite \emph{\'etale} cover of the good open and resolving,
there are semiample holomorphic line bundles $\overline{\mathcal H}_i$
on the smooth compact manifold $P'$ whose specified restrictions are
$\pi^*\mathcal H_i$.

The boundary may be taken SNC.  Along any disk in $P'$ whose punctured
disk lies in the good open, a generator of
$\overline{\mathcal H}_i$ and its inverse have at most powers-of-logarithm
growth in the volume norm~\eqref{periods:volume} and its dual.  The same
is true for all positive powers, including the powers used in the
product-volume identification of Proposition~\ref{spread:product}.
\end{proposition}

\begin{proof}
For a torus-type factor take
$\mathbb V_i=R^1(t_i)_*\Z/\mathrm{tors}$, and for a symplectic-type
factor take $\mathbb V_i=R^2(t_i)_*\Z/\mathrm{tors}$.
The cohomological conditions in Proposition~\ref{spread:product} give
the Hodge numbers in Lemma~\ref{periods:griffiths}.  The fixed ambient
K\"ahler class supplies their flat real polarizations by
Lemma~\ref{periods:lefschetz}; in degree two we use the full
polarization~\eqref{periods:full-h2}.

Apply the finite-level construction simultaneously to these finitely many
variations and to $R^{n_i}(t_i)_*\Z/\mathrm{tors}$ for every $i$.
The latter variations also have flat real polarizations and can be
rationalized by Lemma~\ref{periods:rationalization} for this purpose.
Intersect the resulting finite-index subgroups in the base fundamental
group.  Extending this one finite cover and resolving gives $P'$ with
unipotent local monodromies for all these variations.  We may first
resolve the complement of $P^\circ$ and shrink the good open, so the
final complement is SNC.  Subsequent cup products are taken in the
original real variations, whose monodromies are unchanged by
rationalization.

Lemma~\ref{periods:griffiths} makes the extended Griffiths lines of
$\mathbb V_i$ semiample on this common cover; the same proof uses the
chosen common neat level and requires no additional alteration.
For a torus-type factor, $h^{1,0}=n_i$, and cup product gives the
specified isomorphism
\[
 \bigwedge^{n_i}F^1E_i=\det F^1E_i
       \xrightarrow{\ \sim\ }\mathcal H_i.
\]
It is induced by the pure Hodge morphism
$\bigwedge^{n_i}\mathbb V_i\to R^{n_i}(t_i)_*\R$.
By Lemma~\ref{periods:extensions} it identifies the extended lines.
The determinant calculation above therefore makes
$\overline{\mathcal H}_i$ semiample.

For a symplectic-type factor, write $n_i=2k_i$ and $\ell_i=F^2E_i$.
Equation~\eqref{periods:determinants} makes $\overline\ell_i$ semiample,
and the specified top cup product is
\[
 \ell_i^{\otimes k_i}\xrightarrow{\ \sim\ }\mathcal H_i.
\]
This is induced by the pure morphism
$\mathbb V_i^{\otimes k_i}\to R^{2k_i}(t_i)_*\R$.
Again Lemma~\ref{periods:extensions} identifies the extensions, so
$\overline{\mathcal H}_i\simeq\overline\ell_i^{\otimes k_i}$ is
semiample.  In particular the extended cup maps take generators to
generators, rather than to sections vanishing at the boundary.

We choose $\overline{\mathcal H}_i$ throughout to be the Schmid highest
line in the original top cohomological variation.  Its norm is the
geometric volume norm~\eqref{periods:volume}.  The last assertion is
therefore exactly Lemma~\ref{periods:extensions}, including its pullback
statement for disks.  This proves both the specified holomorphic
extension and the required metric control.
\end{proof}

An absent factor, or a zero-dimensional factor retained by convention,
contributes the trivial line with constant norm.  Thus an empty collection
of period factors requires no modification and no extra hypothesis.
